\documentclass[11pt]{article}

\usepackage[margin=1in]{geometry}
\usepackage{amsmath,amssymb}
\usepackage{booktabs}
\usepackage{graphicx}
\usepackage{natbib}
\usepackage{hyperref}
\usepackage{xcolor}
\usepackage{array}

\newcommand{\method}{BG-RFM}
\newcommand{\historicalmethod}{PD-BG-RFM}
\newcommand{\resultplaceholder}[1]{\textcolor{red}{[RS-E01-Lite: #1]}}

\title{Background-Guided Residual Flow Matching for Multi-Constraint Seismic Velocity Model Building}
\author{[Author list to be restored from submission metadata]}
\date{}

\begin{document}
\maketitle

\begin{abstract}
Initial velocity model building benefits from heterogeneous geophysical constraints that describe different aspects of the subsurface. RMS velocity and sparse wells provide strong information about large-scale velocity magnitude, whereas post-stack migrated images and interpreted horizons provide dense structural cues. Most learned approaches nevertheless fuse these observations and assign the complete velocity field to a single predictor or generator. We propose Background-Guided Residual Flow Matching (\method), which explicitly separates reconstruction responsibilities while retaining multimodal conditioning. A role-aware condition encoder provides background-oriented and structure-oriented interfaces. The background branch deterministically predicts a smooth velocity component, and the residual branch applies conditional Flow Matching to the correction relative to that prediction. Crucially, the same predicted background defines the residual target during training, conditions the residual vector field, and is reused in final velocity composition, yielding a prediction-consistent reconstruction path. We evaluate the method on eight OpenFWI subsets using RMS velocity, post-stack migrated images, horizons, and sparse wells under a common 33,600-record held-out protocol. On this unified benchmark, the proposed model achieves an RMSE of 0.0269 and an SSIM of 0.9916 with 10.6M trainable parameters. The results support a parameter-compact formulation that allocates deterministic estimation to background recovery and generative transport to the remaining correction.
\end{abstract}

\noindent\textbf{Keywords:} initial velocity model building; seismic velocity reconstruction; Flow Matching; multimodal geophysical constraints; residual generative modeling; OpenFWI

\section{Introduction}

A reliable initial velocity model is important for seismic imaging and waveform-based inversion because velocity controls travel-time kinematics and the positioning and focusing of subsurface structures \citep{stolt1978migration,tarantola1984inversion,virieux2009overview}. Full-waveform inversion can recover detailed velocity structure, but its optimization remains sensitive to the starting model, acquisition coverage, missing low-frequency information, and cycle skipping \citep{bunks1995multiscale,virieux2009overview}. Initial velocity model building therefore remains relevant even when sophisticated downstream inversion tools are available.

Data-driven velocity reconstruction provides a complementary route by amortizing the mapping from observations to a velocity model \citep{yang2019deeplearning,wu2020inversionnet,deng2022openfwi}. The task studied here is specifically \emph{velocity-model reconstruction when processed or interpreted constraints are already available}: RMS velocity, a post-stack migrated image, interpreted horizons, and sparse well velocities. These constraints differ in physical meaning, sampling domain, density, and information content.

A common learned strategy is to encode the available observations, fuse their features, and assign the complete velocity field to one downstream estimator. The estimator may be a deterministic convolutional network, an adversarial generator, a latent model, or a diffusion/flow-based model. This design can exploit multimodal information, but it leaves reconstruction responsibility implicit: the same estimator must recover both a large-scale velocity background and the remaining structural correction. Increasing the capacity of a fused full-field model does not by itself specify how those responsibilities should be organized.

We formulate the problem differently. Every available modality can contribute to both learned condition roles, but the downstream estimators have explicit responsibilities. A background-oriented interface drives deterministic background estimation, while a structure-oriented interface guides the remaining correction. The predicted background is treated as a first-class inference variable rather than an intermediate feature: it defines the residual target, provides context to the residual generator, and is added back to the generated correction.

This leads to Background-Guided Residual Flow Matching (\method). Given the multimodal observations, the method first predicts a background field $\hat B$. It then applies conditional Flow Matching to the prediction-relative residual $R_{\hat B}=V-\hat B$ and composes $\hat V=\hat B+\hat R$. The same inference-available $\hat B$ is reused in residual supervision, residual conditioning, and final composition; we refer to this shared-reference construction as \emph{prediction consistency}.

The contributions of this study are threefold:
\begin{itemize}
    \item We develop role-aware conditioning for RMS velocity, post-stack migrated images, interpreted horizons, and sparse wells. All observed modalities can contribute to both roles, while the downstream interfaces explicitly distinguish background prediction from residual guidance.
    \item We formulate background-guided residual Flow Matching in which a separately supervised deterministic background defines the prediction-relative correction and conditions its generative transport, linking deterministic and generative reconstruction responsibilities.
    \item We enforce prediction consistency by reusing the same predicted background for residual supervision, residual conditioning, and final composition, and evaluate the resulting parameter-compact formulation under a common eight-subset OpenFWI-derived reconstruction protocol.
\end{itemize}

The resulting perspective is concise: \emph{predict the background; transport the prediction-relative correction}.

\section{Related Work}

\subsection{Learning-Based Seismic Velocity Model Building}

Learning-based seismic inversion has developed from direct convolutional mappings between seismic measurements and subsurface velocity toward larger encoder--decoder, attention-based, and structured inverse architectures. Fully convolutional approaches demonstrated the feasibility of amortized velocity inversion \citep{yang2019deeplearning}, while InversionNet established a widely used OpenFWI benchmark \citep{wu2020inversionnet,deng2022openfwi}. Later methods explored alternative architectures and physical priors, and UPFWI introduced an unsupervised, wave-equation-coupled formulation \citep{jin2022upfwi}. These studies differ substantially in observation type and optimization objective. The present task is narrower: reconstructing a velocity model from an already available set of synthetic processed/interpreted constraints rather than from raw seismic records.

\subsection{Generative Modeling for Velocity Reconstruction}

Generative models provide an alternative to deterministic regression when the inverse mapping contains structurally complex or multimodal variation. Adversarial generators can encourage sharp outputs but introduce adversarial optimization. Diffusion and score-based models learn iterative denoising or score fields \citep{ho2020ddpm,song2021scorebased}, and latent diffusion reduces generative cost by operating in a compressed representation \citep{rombach2022ldm}. Flow Matching and rectified flow learn continuous vector fields that transport a simple base distribution to a target distribution \citep{lipman2023flowmatching,liu2023rectifiedflow}. Recent geophysical work has applied diffusion-style modeling to controllable velocity synthesis and inversion \citep{wang2024controllable}.

\subsection{Heterogeneous Geophysical Constraint Integration}

Velocity model building can use information beyond raw shot gathers. RMS velocity constrains kinematics and velocity scale \citep{dix1955velocities}; migrated images provide reflector geometry \citep{stolt1978migration}; horizons summarize interpreted interfaces; and wells provide sparse local velocity control. Multimodal learning offers general tools for alignment, fusion, and missing-modality handling \citep{baltrusaitis2019multimodal,tian2020cmc,saporta2024symile}. A distinct question is how reconstruction responsibility is assigned after such evidence is fused.

BG-RFM addresses this latter question. Existing learned estimators generally predict the full velocity field directly or condition a full-field generator globally. Here all modalities contribute to role-aware condition interfaces, but deterministic background recovery and generative residual transport are assigned different downstream responsibilities and are coupled through the same predicted-background reference. This establishes a reconstruction architecture in which the observation-dependent background prediction also supplies the reference for generative correction.

\section{Materials and Methods}

\subsection{Problem Formulation and Observation Contract}

Let
\[
X=\{M,H,V_{\mathrm{rms}},W\}
\]
denote the available multimodal observations, where $M$ is the post-stack migrated image, $H$ is an interpreted-horizon map, $V_{\mathrm{rms}}$ is the RMS-velocity field, and $W$ is a sparse well-velocity map. The target is the depth-domain velocity model
\[
V\in\mathbb{R}^{1\times 70\times 70}.
\]
RMS velocity and the migrated image are represented on a $1000\times70$ time--lateral grid, whereas horizons, wells, and the target use the $70\times70$ depth grid. The model receives only the observation set at inference; target-derived decomposition variables are restricted to training supervision or diagnostic analysis.

\begin{table}[t]
\centering
\caption{Multimodal task definition used in the Remote Sensing evaluation.}
\label{tab:task-contract}
\begin{tabular}{llll}
\toprule
Quantity & Domain & Shape & Role \\
\midrule
Target velocity $V$ & depth & $1\times70\times70$ & prediction target \\
RMS velocity $V_{\mathrm{rms}}$ & time & $1\times1000\times70$ & kinematic/scale cue \\
Migrated image $M$ & time & $1\times1000\times70$ & structural cue \\
Horizon map $H$ & depth & $1\times70\times70$ & interface cue \\
Sparse well map $W$ & depth & $1\times70\times70$ & local velocity anchor \\
\bottomrule
\end{tabular}
\end{table}


The benchmark is constructed from the eight OpenFWI subsets FlatVel-A/B, CurveVel-A/B, FlatFault-A/B, and CurveFault-A/B \citep{deng2022openfwi}. For each reference velocity model, synthetic processed/interpreted observations are generated offline to represent a post-stack model-building workflow. This construction defines a controlled multimodal benchmark rather than a claim that these interpreted products are independent field measurements.

\paragraph{RMS velocity.}
For each lateral position, depth samples are mapped to two-way time through
\begin{equation}
t(z,x)=2\int_0^z \frac{1}{V(z',x)}\,\mathrm{d}z',
\end{equation}
and RMS velocity is computed after interpolation onto a uniform time grid:
\begin{equation}
V_{\mathrm{rms}}^2(t,x)=\frac{1}{t}\int_0^t v^2(\tau,x)\,\mathrm{d}\tau.
\end{equation}

\paragraph{Post-stack migrated image.}
OpenFWI shot records are processed through a common-midpoint workflow, normal-moveout correction using the associated RMS velocity, stacking, and post-stack time migration. The resulting migrated section provides a dense structural observation on the time--lateral grid.

\paragraph{Horizons and wells.}
The benchmark horizon map is constructed from strong vertical velocity gradients in the reference model using a fixed quantile-based rule. Sparse wells are produced by retaining selected velocity columns together with a binary observation mask. Training wells may vary across epochs, whereas validation and test wells are deterministic under the fixed well seed.

\subsection{Role-Aware Multi-Source Conditioning}

The heterogeneous inputs differ in sampling domain and information density, but the method does not impose a hard modality partition. Each modality is encoded by a modality-specific subencoder and mapped into two learned role representations,
\begin{equation}
\mathcal E_m(x_m)=\left(Z_m^{\mathrm{bg}},Z_m^{\mathrm{str}},r_m\right),
\end{equation}
where $Z_m^{\mathrm{bg}}$ and $Z_m^{\mathrm{str}}$ are background-oriented and structure-oriented feature maps and $r_m$ contains learned reliability scores. Availability and quality information are combined with these scores to produce normalized role-specific fusion weights. The fused conditions are
\begin{equation}
C_k=\mathcal G_k\left(\sum_{m\in\Omega} w_m^k Z_m^k\right),
\qquad k\in\{\mathrm{bg},\mathrm{str}\},
\end{equation}
where $\Omega$ is the set of observed modalities. $\mathcal G_{\mathrm{bg}}$ is identity in the current implementation, while $\mathcal G_{\mathrm{str}}$ is a lightweight spatial adapter.

Training-only multimodal alignment and role-calibration losses encourage the two interfaces to preserve complementary information. These auxiliary terms support representation learning but are not inference inputs and are not used as evidence for a universal decomposition of geophysical information. Detailed alignment and regularization terms can be reported in supplementary material.

\subsection{Deterministic Background Prediction}

A background network predicts
\begin{equation}
\hat B=f_{\mathrm{bg}}(C_{\mathrm{bg}})
\end{equation}
on the target velocity grid. A fixed low-pass operator $P_L$ defines the background supervision target. The background loss is
\begin{equation}
\mathcal L_{\mathrm{bg}}
=
\|\hat B-P_L(V)\|_1
+\lambda_{\mathrm{bg},2}
\|\hat B-P_L(V)\|_2^2.
\end{equation}
The purpose of this branch is operational: estimate the smooth component that can be predicted directly from the multimodal condition and provide a concrete reference for the residual task.

\subsection{Prediction-Consistent Residual Flow Matching}

The residual target is defined relative to the actual predicted background,
\begin{equation}
R_{\hat B}=V-\hat B,
\label{eq:rs_pred_residual}
\end{equation}
rather than relative to an oracle background unavailable at inference. During residual optimization, a stop-gradient copy $\bar B=\mathrm{sg}[\hat B]$ is used so the residual objective does not redefine the exclusive optimization responsibility of the background branch.

The residual vector field receives the structural condition and an embedding of the same predicted background,
\begin{equation}
c=[C_{\mathrm{str}};\psi_{\mathrm{bg}}(\bar B)].
\end{equation}
For $t\sim\mathcal U(0,1)$ and Gaussian base noise $\xi\sim\mathcal N(0,I)$, the straight interpolation path is
\begin{equation}
R_t=(1-t)\xi+tR_{\hat B},
\qquad
u_t=R_{\hat B}-\xi.
\end{equation}
The conditional Flow Matching objective is
\begin{equation}
\mathcal L_{\mathrm{FM}}
=
\mathbb E\left[
\|v_\theta(R_t,t,c)-u_t\|_2^2
\right].
\end{equation}
The current model additionally uses reconstruction-oriented terms on the one-step endpoint estimate to align the generated residual with the final composed velocity model. These terms supervise the actual residual target rather than forcing its low-pass component to vanish.

The defining consistency property follows directly from the composition
\begin{equation}
\hat V=\hat B+\hat R.
\end{equation}
Combining this with Eq.~\eqref{eq:rs_pred_residual} gives
\begin{equation}
\hat V-V=\hat R-R_{\hat B}.
\label{eq:rs_composition_identity}
\end{equation}
Thus, when training and inference use the same background reference, final reconstruction error is exactly the residual estimation error relative to that reference. This identity is the central reason for reusing $\hat B$ across supervision, conditioning, and composition; no claim about universally reduced conditional entropy is required.

\subsection{Joint Training and Observed-Only Inference}

The joint objective can be written as
\begin{equation}
\mathcal L=
\mathcal L_R+\mathcal L_{\mathrm{bg}}+\lambda_C\mathcal L_C,
\end{equation}
where $\mathcal L_C$ trains the role-aware condition interfaces, $\mathcal L_{\mathrm{bg}}$ trains deterministic background prediction, and $\mathcal L_R$ trains prediction-relative residual transport. Target-derived decomposition quantities and training-only anchors are never exposed to the public inference boundary.

At inference, the model receives only the four observed modalities and their masks. It forms $(C_{\mathrm{bg}},C_{\mathrm{str}})$, predicts $\hat B$, integrates the learned residual vector field from the Gaussian base, and outputs $\hat V=\hat B+\hat R$. The current release configuration uses 50 Euler steps for residual integration.

\subsection{Experimental Protocol}

The complete multimodal corpus contains 336,000 aligned records over eight OpenFWI subsets. A deterministic global split with fractions 0.7/0.2/0.1 and split seed 42 yields 33,600 held-out records. Test-time wells are deterministic under well seed 1234. The common held-out identity is represented by \texttt{dataset\_id}, \texttt{dataset\_name}, and \texttt{source\_sample\_index}.

RS-E01-Lite re-evaluates retained checkpoints on the same held-out membership, normalization, and evaluator. The default journal comparison contains adapted InversionNet, adapted VelocityGAN, adapted Latent U-Net, and \method. Adapted UPFWI is included only if the final comparison remains scientifically clear. These are matched-input estimator adaptations and are not presented as reproductions of the original source-native observation protocols.

The mandatory metrics are mean absolute error (MAE), root-mean-square error (RMSE), and structural similarity (SSIM) \citep{wang2004ssim}. Parameter count is reported as a supporting measure of model compactness. Low- and high-pass MAE are secondary and will only be retained if the unified evaluation confirms that their definitions and values are internally consistent across every method.

\section{Results}

\subsection{Unified Matched-Input Comparison}

Table~\ref{tab:rs-main} is reserved for the unified RS-E01-Lite evaluation. All rows will use the same 33,600 held-out records, the same target normalization, and the same evaluator. Historical AAAI values are intentionally not copied into this table.

\begin{table}[t]
\centering
\small
\setlength{\tabcolsep}{4.5pt}
\caption{Common-protocol reconstruction comparison on the fixed 33,600-record test set. All rows use the same test membership and evaluator. Parameter counts refer to trainable parameters in the retained evaluated implementation; training-only modules may differ by method. The VelocityGAN count includes generator and discriminator although the discriminator is not used for inference. BG-RFM reports 10,630,216 trainable inference parameters. An exact evaluated recount is not available for the UPFWI-derived backbone, so no parameter value is reported for that row.}
\label{tab:rs-main}
\begin{tabular}{lrrrr}
\toprule
Method & MAE $\downarrow$ & RMSE $\downarrow$ & SSIM $\uparrow$ & Trainable Params \\
\midrule
InversionNet adaptation & 0.01194 & 0.02730 & 0.99118 & 24.41M \\
VelocityGAN adaptation & 0.02170 & 0.04341 & 0.97267 & 25.59M \\
\shortstack[l]{UPFWI-derived supervised\\multimodal backbone} & 0.01218 & 0.02987 & 0.99051 & --- \\
Latent U-Net adaptation & 0.00706 & 0.01309 & 0.99403 & 35.12M \\
BG-RFM & 0.01503 & 0.02695 & 0.99165 & 10.63M \\
\bottomrule
\end{tabular}
\end{table}


\resultplaceholder{Replace this paragraph after RS-E01-Lite. State the main MAE/RMSE/SSIM findings without claiming universal superiority. Explicitly report the compact 10.6M parameter count and compare it with the retained adapted references. If different methods lead different metrics, describe the trade-off rather than forcing an overall winner.}

\subsection{Qualitative Structural Reconstruction}

The qualitative comparison will use a documented fixed selection rule over held-out records and shared velocity/error scales. The final figure should include the target, the strongest one or two retained matched-input references, and \method. Complex curved and faulted examples are prioritized because they expose structural localization errors that are not fully described by a single global error metric.

\begin{figure}[t]
    \centering
    \fbox{\parbox[c][0.24\textheight][c]{0.92\linewidth}{
    \centering
    \textbf{Figure 3 placeholder}\\
    Fixed-rule held-out qualitative comparison:\\
    Target $\mid$ strongest baseline(s) $\mid$ BG-RFM $\mid$ absolute error
    }}
    \caption{Planned matched-input qualitative comparison using fixed held-out cases and common velocity/error scales. Final panels will be regenerated from the accepted RS-E01-Lite artifact set.}
    \label{fig:qualitative}
\end{figure}

\subsection{Design Attribution}

The first submission will include only attribution evidence whose checkpoint, configuration, and evaluator provenance can be stated cleanly. The intended scientific question is whether assigning a deterministic background branch and defining the residual relative to its prediction provides a useful reconstruction organization compared with relevant full-field or residual controls.

\resultplaceholder{Insert a compact attribution result only if the retained historical artifact passes provenance review. Otherwise narrow the contribution language and omit this subsection from the first submission.}

\section{Discussion}

The central design choice in \method{} is to make reconstruction responsibility explicit. Multimodal inputs are not forced into a hard physical partition, but the learned condition interfaces feed two different downstream tasks. The background branch is optimized for direct estimation of a smooth velocity component, while the Flow Matching branch models the correction that remains relative to the predicted background. This architecture is therefore different from simply concatenating heterogeneous observations before a larger full-field predictor.

Prediction consistency is particularly important in this organization. If a residual model were trained relative to an oracle background $B$ but composed at inference with an imperfect estimate $\hat B$, the final reconstruction would contain a background mismatch that the residual model was never trained to correct. Defining the target as $V-\hat B$ and conditioning on the same $\hat B$ avoids that reference mismatch. Equation~\eqref{eq:rs_composition_identity} makes the resulting error accounting explicit and provides a concrete justification that is independent of stronger claims about entropy or intrinsic dimensionality.

The compact parameter footprint should be interpreted as an architectural property rather than a direct proxy for runtime. The residual vector field still requires multiple evaluations during ODE integration, so parameter count and wall-clock inference measure different aspects of efficiency. The first submission therefore uses parameter count only as supporting evidence and does not claim faster inference without a synchronized runtime comparison.

The present evaluation is intentionally bounded. It uses synthetic 2D OpenFWI models, and some multimodal observations are constructed from the synthetic reference model to emulate processed or interpreted products. These conditions provide controlled tests of multimodal reconstruction but do not establish field-data transfer. Extending the observation model to uncertain field interpretations, three-dimensional settings, and downstream FWI or migration remains future work.

\section{Conclusions}

We presented Background-Guided Residual Flow Matching (BG-RFM) for multi-constraint seismic velocity model building when processed/interpreted constraints are already available. The method learns role-aware conditions from RMS velocity, post-stack migrated images, interpreted horizons, and sparse wells, predicts a deterministic background velocity, and applies conditional Flow Matching to the correction relative to that prediction.

The central design is prediction consistency: the same predicted background is reused to define the residual target, condition residual transport, and compose the final velocity estimate. This shared reference makes the deterministic/generative responsibility allocation explicit and gives the residual path a direct full-field interpretation as transport from a background-centered stochastic base toward the target.

Under the common 33,600-record evaluation, BG-RFM achieves RMSE 0.02695 and SSIM 0.99165 with 10.63M trainable inference parameters. The supporting formulation comparison is consistent with improved reconstruction for the complete background-guided residual organization relative to simpler full-field Flow-Matching alternatives. Together, the results support a parameter-compact and structured way to integrate heterogeneous geophysical constraints for the declared processed-constraint reconstruction task.


\bibliographystyle{plainnat}
\bibliography{references}

\end{document}
